\documentclass[10pt,a4paper]{amsart}
\usepackage[margin=19mm]{geometry}
\usepackage{amsmath,amssymb,mathtools}
\usepackage[scr=rsfs,cal=euler]{mathalfa}
\usepackage{xfrac}
\usepackage[unicode,hidelinks,bookmarks=false]{hyperref}
\newcommand{\ie}{i.e.\ }
\newcommand{\cf}{cf.\ }
\newcommand{\resp}{respectively\ }
\newcommand{\Subsection}{\subsection}
\providecommand{\nobreakdash}{\nobreak\hbox{-}}
\setlength{\emergencystretch}{2em}
\multlinegap0pt
\usepackage{bm}
\usepackage{bbm}
\usepackage{dsfont}
\usepackage{xspace}
\usepackage{cancel}
\usepackage{mathtools}
\usepackage{amsthm}
\makeatletter
\@ifundefined{theorem}{\newtheorem{theorem}{Theorem}}{}
\@ifundefined{lemma}{\newtheorem{lemma}[theorem]{Lemma}}{}
\makeatother
\DeclareMathOperator{\cum}{cum}
\DeclareMathOperator{\tr}{tr}
\title[Counterexamples to the Mao-Wu-Xu spectral conjecture and a t]{Counterexamples to the Mao-Wu-Xu spectral conjecture and a three-scale boundary for spherical random graphs}
\author{Congyi Luo}
\date{Source: arXiv:2609.11662v2 (CC BY 4.0)}
\begin{document}
\maketitle

\noindent\emph{Source and licence.} Counterexamples to the Mao-Wu-Xu spectral conjecture and a three-scale boundary for spherical random graphs. Version arXiv:2609.11662v2 (CC BY 4.0), distributed under the Creative Commons Attribution 4.0 International licence, as recorded in the arXiv licence metadata for this exact version and shown on the abstract page https://arxiv.org/abs/2609.11662v2. Full rights notice: licenses/arxiv-2609.11662-CC-BY-4.0.txt.\par\smallskip

\noindent\textit{Editorial scope.} Counted unit: the Lemma whose source label is \texttt{lem:core-sum} in the article (source0.tex lines 837--846, source label \texttt{lem:core-sum}), delivered together with the article's own complete original proof of it (source0.tex lines 847--939, one \texttt{proof} environment of 93 lines). No mathematics is written, repaired or re-derived: statement and proof are the article's own text line for line. Required context retained uncounted: lem:moments (lines 341--353); eq:h-cube (lines 350--352); eq:cumulant-def (lines 786--790); lem:cut (lines 809--813). Every definition, display and label of those lines is retained unchanged. Omitted from this package and not redistributed: 1--340, 353--785, 791--808, 814--836, 940--1621. Nothing inside a retained excerpt is omitted or altered: every retained body is delivered exactly as the article's lines are written, so no omission receipt of any kind accompanies this package. Citations: the retained excerpts carry no citation command at all, so no bibliography entry of the article is transcribed and no reference is redistributed. Reference adaptation: none was needed and none was made; every reference inside the delivered unit resolves inside it, so no unresolved reference or citation is produced. Printed slips kept literal: no printed word, symbol, hypothesis or number of the counted statement or of its proof was altered. Layout and class: the article's own class is \texttt{amsart} with packages including etoolbox, geometry, amsmath, amssymb, amsthm, booktabs, hyperref; the wrapper is the lane's pinned \texttt{amsart} editorial wrapper at 10pt on A4, which restates the theorem-like environments the retained text needs and adds the article's own macro definitions that the retained text needs (\texttt{\textbackslash cum}, \texttt{\textbackslash tr}), with extra wrapper packages (bm, bbm, dsfont, xspace, cancel, mathtools). The delivered numbering is the unit's own. Assets: no retained span contains a figure environment, a graphics-inclusion command, a PDF-page inclusion, a table or a TikZ picture; no image file is copied into this package, the package declares no asset dependency and no image file is redistributed with it. Licence: version arXiv:2609.11662v2 of this article is distributed under the Creative Commons Attribution 4.0 International licence, as recorded in the arXiv licence metadata for this exact version and shown on the abstract page https://arxiv.org/abs/2609.11662v2. A full rights notice is delivered at licenses/arxiv-2609.11662-CC-BY-4.0.txt. Characters: every retained excerpt is pure ASCII, so the delivered engine, XeLaTeX with its default Latin Modern text font, reports no missing character.
\begin{lemma}\label{lem:moments}
For each fixed positive integer $q$, there is a constant $C_{L,q}$ depending only on $L,q$
such that every real polynomial $p$ of degree at most $L$ satisfies
\[
 \mathbb E|p(\langle X,Y\rangle)|^{2q}
 \le C_{L,q}\bigl(\mathbb E p(\langle X,Y\rangle)^2\bigr)^q,
\]
where $X,Y$ are independent with law $\mu_d$, uniformly for all $d\ge2$.
Moreover, for the bounded centered kernel $k$ above,
\begin{equation}\label{eq:h-cube}
 h^3\le C_{L,b}\tau.
\end{equation}
\end{lemma}

\begin{equation}\label{eq:cumulant-def}
 \cum(Z_1,\ldots,Z_m)
 =\sum_{\pi\in\mathcal P([m])}(-1)^{|\pi|-1}(|\pi|-1)!
      \prod_{B\in\pi}\mathbb E\prod_{j\in B}Z_j.
\end{equation}

\begin{lemma}\label{lem:cut}
If the underlying simple graph of $H$ is disconnected or has a cut vertex, then $C_H=0$.
Also, $C_H=0$ if $H$ has a vertex of degree one, counting multiplicity.
Throughout, degrees count edge multiplicities, and a cut vertex is a vertex whose deletion increases the number of connected components.
\end{lemma}

\begin{lemma}\label{lem:core-sum}
Fix an integer $t_0\ge1$ and a real number $w\ge1$. For $A_4$ sufficiently small,
\begin{equation}\label{eq:core-sum}
 \sum_{\substack{H\text{ connected}\\m\ge3,\ t\le t_0}}
       w^{m-1}\frac{m\,C_H^2}{\prod_fm_f!}
 \le C_{L,b,t_0,w}\bigl(A_3+A_4^{1/6}\bigr).
\end{equation}
The sum is over distinct labeled multigraphs with vertices in $[n]$ and no isolated vertices,
without an additional ordering of the edge occurrences.
\end{lemma}

\begin{proof}
By Lemma~\ref{lem:cut}, only supports without cut vertices need be considered.
First separate two cases with no nontrivial core.

If the support has only two vertices, then $m=t+1\le t_0+1$.
By \eqref{eq:cumulant-def} and Lemma~\ref{lem:moments},
$|C_H|\le C_{L,t_0}h^{m/2}$.
Indeed, each moment factor of size $r$ is at most $C_{L,t_0}h^{r/2}$,
and the block sizes sum to $m$.
Since $h\le b^2$ and $3\le m\le t_0+1$, we have $h^m\le C_{b,t_0}h^3$, so the total contribution is at most
$C n^2h^3\le C n^2\tau=C\sqrt{A_4}$.

If all degrees equal two and there are at least three vertices,
then $H$ is a simple cycle. Every nontrivial partition in the cumulant
puts two adjacent edges in different blocks; one moment factor then has a degree-one vertex and vanishes.
Hence $C_H=\tr T^m$. Triangles contribute $O_w(A_3)$ in total.
For $m\ge4$, the finite spectral sum gives
\[
 |\tr T^m|\le \tau\|T\|_{\mathrm{op}}^{m-4},
 \qquad n^m|\tr T^m|^2\le A_4\rho^{m-4}.
\]
There are at most $n^m$ labeled $m$-cycles, so these cycles contribute at most
\[
 C_w A_4\sum_{m\ge4}m(w\rho)^{m-4}\le C_wA_4
\]
provided $w\rho\le1/2$.

Now suppose the support has at least three vertices and is not one of the cycles above.
Suppress the degree-two vertices successively. Since the support has no cut vertex,
each such vertex is incident to two distinct edges of multiplicity one: if it were joined only to a single neighbor by a repeated edge,
that neighbor would be a cut vertex. Call the remaining vertices core vertices, and denote their number by $v_0$.
There are $R_0$ direct edges between core vertices, counted with multiplicity, and $p$ further paths
of length at least two. Let their lengths be $r_1,\ldots,r_p$,
and set $z_i=r_i-2$, $z=\sum_i z_i$, and $B=R_0+p$.
Then
\begin{equation}\label{eq:core-counts}
 v=v_0+p+z,\qquad m=R_0+2p+z,\qquad
 t=B-v_0+1,\qquad 3v_0\le2B.
\end{equation}
Thus $v_0\le2(t_0-1)$ and $B\le3(t_0-1)$.
This case is empty when $t_0=1$.

Each suppressed path has distinct endpoints. Otherwise it meets the rest of the graph at only one core vertex,
which would be a cut vertex; if there were no other vertices, the original graph would be a cycle.
Adding a repeated edge to a cycle creates two distinct core endpoints, so this gives no exception.
Thus every convolution kernel below is evaluated at two independent spherical points.

In every nonzero term of \eqref{eq:cumulant-def}, all edges of a given path belong to one block;
otherwise a degree-one vertex appears at a separation point. Each long path may therefore be treated as one unit,
and each occurrence of a direct edge as another unit.
Nonzero partitions are partitions of these $B$ units, so their number and the coefficients $(|\pi|-1)!$ depend only on $t_0$.
Within each moment factor, integrate the internal vertices of a path of length $r$ to replace it by the integral kernel of $T^r$.
This is still an inner-product polynomial of degree at most $L$, and
\[
 \|T^r\|_{\mathrm{HS}}^2
   =\sum_\ell M_{\ell,d}(a_\ell/M_{\ell,d})^{2r}
   \le \tau\|T\|_{\mathrm{op}}^{2r-4},\qquad r\ge2.
\]
Apply H\"older's inequality to the finitely many factors on the core and then Lemma~\ref{lem:moments}
to obtain
\begin{equation}\label{eq:core-moment}
 C_H^2\le C_{L,t_0}h^{R_0}\tau^p\|T\|_{\mathrm{op}}^{2z}.
\end{equation}
H\"older's inequality does not require independence of the factors. Their individual endpoints are distinct,
which ensures that their squared $L^2$ norms equal the corresponding squared Hilbert--Schmidt norms.

For a fixed core and length vector, first number the core vertices and orient the paths,
then fill in the labels of the core vertices and internal path vertices in order.
The bounded core size gives at most $C_{t_0}n^v$ such encodings.
Every graph has at least one encoding; allowing repeated labels only increases the upper bound.
No additional factor of $m!$ or $v!$ is needed.
The weight satisfies $m/\prod_fm_f!\le m$.
When $A_4\le1$, substituting \eqref{eq:h-cube} into \eqref{eq:core-moment} gives
\begin{align}
 n^vC_H^2
 &\le C_{L,b,t_0}
  n^{\,v_0-2R_0/3-p} A_4^{R_0/6+p/2}\rho^z\notag\\
 &\le C_{L,b,t_0}A_4^{1/6}\rho^z.\label{eq:core-small}
\end{align}
For the second inequality, $v_0\le2(R_0+p)/3$, so the exponent of $n$ is at most $-p/3$;
also $R_0+p\ge1$, so the exponent of $A_4$ is at least $1/6$.

Finally, sum over path lengths. With $a=R_0+2p$ and $x=w\rho$,
\[
 \sum_{z_1,\ldots,z_p\ge0}(a+\textstyle\sum_i z_i)
       w^{a+\sum_i z_i-1}\rho^{\sum_i z_i}
 =w^{a-1}\left\{\frac{a}{(1-x)^p}
             +\frac{px}{(1-x)^{p+1}}\right\}.
\]
This identity follows from the geometric series and its derivative. When $p=0$, the sum has one empty length vector and equals $w^{a-1}a$.
For $w\rho\le1/2$, the right-hand side depends only on $t_0,w$.
There are finitely many core types. Adding the repeated-edge and cycle contributions proves \eqref{eq:core-sum}.
\end{proof}

\end{document}
